\chapter*{不定积分每日一题}

\refer{《积分的方法与技巧》金玉明\ 顾新身\ 毛瑞庭\ /\ 编著}

\exercise{求积分$\int{\ee^{ax}\cos{bx}\d x},\ \int{\ee^{ax}\sin{bx}\d x}$}

\teacher{陆丽芬}

\postdate{2022年12月4日}

\bigresult{
    \formula{\int{\ee^{ax}\cos{bx}\d x}=\ee^{ab}\frac{a\cos{bx}+b\sin{bx}}{a^2+b^2}+C}
    \formula{\int{\ee^{ax}\sin{bx}\d x}=\ee^{ab}\frac{a\cos{bx}-b\sin{bx}}{a^2+b^2}+C}
}

\analysis{
    这是一道非常经典的分部积分题，运用反对幂指三的经验公式可以操作.
    \formula{
        \begin{aligned}
            I&=\int{\ee^{ax}\cos{bx}}=\frac{1}{b}\int{\ee^{ax}\d\sin{bx}}\\
            &=\frac{1}{b}\qty(\ee^{ax}\sin{bx}-a\int{\sin{bx}\ee^{ax}\d x})\\
            &=\frac{1}{b}\qty(\ee^{ax}\sin{bx}+\frac{a}{b}\int{\ee^{ax}\d\cos{bx}})\\
            &=\frac{1}{b}\qty(\ee^{ax}\sin{bx}+\frac{a}{b}\ee^{ax}\cos{bx}-\frac{a^2}{b}\int{\ee^{ax}\cos{bx}\d x})\\
            &=\frac{1}{b}\qty(\ee^{ax}\sin{bx}+\frac{a}{b}\ee^{ax}\cos{bx}-\frac{a^2}{b}I)
        \end{aligned}
    } 
    两次分部积分之后，原本要求的积分又出现了，因此我们可以解方程（不要忘记加上最后的常数$C$哦$\sim$，这种操作确实很容易忘记）：
    \formula{bI=\ee^{ax}\sin{bx}+\frac{a}{b}\ee^{ax}\cos{bx}-\frac{a^2}{b}I}
    \formula{\frac
    {a^2+b^2}{b}I=\ee^{ax}\sin{bx}+\frac{a}{b}\ee^{ax}\cos{bx}}
    \formula{I=\frac{b\ee^{ax}\sin{bx}+a\ee^{ax}\cos{bx}}{a^2+b^2}+C}
    而对于另外一个积分$\int{\ee^{ax}\sin{bx}\d x}$，也是同样的操作，就不再演示了，同学们可以自行操作.但其实我们在求第一个积分时，第一次分部之后，出现了$\int{\ee^{ax}\sin{bx}\d x}$，那么根据对称，我们想，如果对$\int{\ee^{ax}\sin{bx}\d x}$也做一次分部积分，应该也会出现$\int{\ee^{ax}\cos{bx}\d x}$，那么就相当于两个方程两个未知数.也可以解出答案，实操一下：


    这里我们用反对幂三指（指数和三角地位差不多，多演示一个），得到
    \formula{\int{\ee^{ax}\cos{bx}\d x}=\frac{1}{a}\ee^{ax}\cos{bx}-\int{\frac{1}{a}\ee^{ax}\qty(-b\sin{bx}\d x)}}
    整理一下可得
    \formula{\int{\ee^{ax}\cos{bx}\d x}=\frac{1}{a}\ee^{ax}\cos{bx}+\frac{b}{a}\int{\ee^{ax}\sin{bx}\d x}}
    同理有
    \formula{\int{\ee^{ax}\sin{bx}\d x}=\frac{1}{a}\ee^{ax}\sin{bx}-\frac{b}{a}\int{\ee^{ax}\cos{bx}\d x}}
    那么把$\int{\ee^{ax}\cos{bx}\d x},\ \int{\ee^{ax}\sin{bx}\d x}$看成两个未知量，即可解得方程
    $$\int{\ee^{ax}\cos{bx}\d x}=\ee^{ab}\frac{a\cos{bx}+b\sin{bx}}{a^2+b^2}+C$$
    $$\int{\ee^{ax}\sin{bx}\d x}=\ee^{ab}\frac{a\cos{bx}-b\sin{bx}}{a^2+b^2}+C$$
    在有三角函数分部积分法中，写出另外一个相对的，然后解方程操作是比较常见的.或者连续使用分部积分法，然后出现原本的式子再解方程，这也是很常见的操作.
}